\frfilename{mt213.tex}
\versiondate{13.9.13}

\copyrightdate{2000}
\def\chaptername{Taksonomi ruang ukur}
\def\sectionname{Ruang semihingga, ditentukan secara lokal, dan dapat dilokalkan}

\newsection{213}

Pada bagian ini saya menghimpun berbagai fakta berguna tentang jenis-jenis
ruang ukur tersebut.   Saya mulai dengan sifat-sifat khas ruang semihingga
(213A-213B), lalu beralih ke ruang lengkap yang ditentukan secara lokal
(213C) dan konsep `versi c.l.d.' (213D-213H),
%213D 213E 213F 213G 213H
yakni metode paling ampuh di antara metode yang selalu tersedia untuk
memodifikasi suatu ruang ukur agar sifat-sifatnya lebih baik.   Saya kemudian
membahas secara singkat `himpunan terabaikan yang ditentukan secara lokal'
(213I-213L),
%213I 213J 213K 213L
dan selubung terukur (213L-213M), lalu menutup bagian ini dengan hasil-hasil
tentang ruang yang dapat dilokalkan (213N) dan ruang yang dapat dilokalkan
secara ketat (213O).

\leader{213A}{Lema} Misalkan $(X,\Sigma,\mu)$ suatu ruang ukur semihingga.
Maka

\Centerline{$\mu E
=\sup\{\mu F:F\in\Sigma,\,F\subseteq E,\,\mu F<\infty\}$}

\noindent untuk setiap $E\in\Sigma$.

\proof{ Tetapkan
$c=\sup\{\mu F:F\in\Sigma,\,F\subseteq E,\,\mu F<\infty\}$.   Jelas
$c\le\mu E$, sehingga jika
$c=\infty$ tidak ada lagi yang perlu dibuktikan.   Jika $c<\infty$, ambil
barisan $\sequencen{F_n}$ yang terdiri atas subhimpunan terukur dari $E$
dengan ukuran berhingga, sedemikian sehingga
$\lim_{n\to\infty}\mu F_n=c$;  tetapkan $F=\bigcup_{n\in\Bbb N}F_n$.
Untuk setiap $n\in\Bbb N$, $\bigcup_{k\le n}F_k$ merupakan himpunan terukur
berukuran berhingga yang termuat dalam $E$, sehingga
$\mu(\bigcup_{k\le n}F_k)\le c$,
dan

\Centerline{$\mu F=\lim_{n\to\infty}\mu(\bigcup_{k\le n}F_k)\le c$.}

\noindent Di pihak lain,


\Centerline{$\mu F\ge\sup_{n\in\Bbb N}\mu F_n\ge c$,}

\noindent jadi $\mu F=c$.

Jika $F'$ merupakan subhimpunan terukur dari $E\setminus F$ dan
$\mu F'<\infty$, maka $F\cup F'$ berukuran berhingga dan termuat dalam $E$,
sehingga ukurannya paling besar $c=\mu F$;  akibatnya $\mu F'=0$.   Ini
berarti bahwa $\mu(E\setminus F)$ tidak mungkin tak berhingga, sebab jika
demikian, kesemihinggaan $(X,\Sigma,\mu)$ mengharuskan adanya suatu himpunan
terukur dengan ukuran berhingga tak nol di dalamnya.   Jadi
$E\setminus F$ berukuran berhingga, dan dengan demikian sebenarnya
terabaikan;  maka $\mu E=c$, sebagaimana dinyatakan.
}%end of proof of 213A

\vleader{72pt}{213B}{Proposisi} Misalkan $(X,\Sigma,\mu)$ suatu ruang ukur
semihingga.   Misalkan $f$ suatu fungsi terukur secara virtual terhadap $\mu$,
bernilai $[0,\infty]$, dan terdefinisi hampir di mana-mana pada $X$.   Maka

$$\eqalign{\int f
&=\sup\{\int g:g\text{ adalah fungsi sederhana},\,g\leae f\}\cr
&=\sup_{F\in\Sigma,\mu F<\infty}\int_Ff}$$

\noindent di $[0,\infty]$.

\proof{{\bf (a)} Pada sebarang ruang ukur $(X,\Sigma,\mu)$, suatu fungsi
bernilai $[0,\infty]$ yang terdefinisi pada subhimpunan $X$ dapat
diintegralkan jika dan hanya jika terdapat himpunan koterabaikan $E$ sehingga

\inset{($\alpha$) $E\subseteq\dom f$ dan $f\restr E$ adalah terukur,}

\inset{($\beta$)
$\sup\{\int g:g$ adalah fungsi sederhana, $g\leae f\}$ berhingga,}

\inset{($\gamma$) untuk setiap $\epsilon>0$,
$\{x:x\in E,\,f(x)\ge\epsilon\}$ berukuran berhingga,}

\inset{($\delta$) $f$ bernilai berhingga hampir di mana-mana}

\noindent (lihat 122Ja, 133B).   Namun, jika $\mu$ semihingga, ($\gamma$)
dan ($\delta$) mengikuti dari syarat-syarat lainnya.   \Prf\ Ambil
$\epsilon>0$. Tetapkan

\Centerline{$E_{\epsilon}=\{x:x\in E,\,f(x)\ge\epsilon\}$,}

\Centerline{$c
=\sup\{\int g:g$ adalah fungsi sederhana, $g\leae f\}$;}

\noindent menurut asumsi, $c$ berhingga.   Jika
$F\subseteq E_{\epsilon}$ adalah terukur dan $\mu F<\infty$, maka
$\epsilon\chi F$ adalah
fungsi sederhana dan $\epsilon\chi F\leae f$, jadi

\Centerline{$\epsilon\mu F=\int\epsilon\chi F\le c$,
\quad$\mu F\le\Bover{c}{\epsilon}$.}

\noindent Karena $F$ sebarang, 213A menunjukkan bahwa
$\mu E_{\epsilon}\le\bover{c}{\epsilon}$ berhingga.   Karena $\epsilon$
sebarang, ($\gamma$) terpenuhi.

Untuk ($\delta$), jika $F=\{x:x\in E,\,f(x)=\infty\}$, maka $\mu F$
berhingga (oleh ($\gamma$)) dan $n\chi F\leae f$, sehingga $n\mu F\le c$
untuk setiap $n\in\Bbb N$;  jadi $\mu F=0$.\ \Qed

\medskip

{\bf (b)} Sekarang misalkan $f:D\to[0,\infty]$ terukur secara virtual
terhadap $\mu$, dengan $D\subseteq X$ koterabaikan, sehingga
$\int f$ terdefinisi dalam $[0,\infty]$ (135F).   Menurut (a),

\leaveitout{$$\eqalignno{\int f
&=\sup_{g\text{ sederhana},g\le f\text{ hampir di mana-mana}}\int g\cr
\displaycause{jika salah satunya berhingga, dan dengan demikian juga jika
salah satunya tak berhingga}
&=\sup_{g\text{ sederhana},g\le f\text{ hampir di mana-mana},\mu F<\infty}\int_Fg
\le\sup_{\mu F<\infty}\int_Ff
\le\int f,\cr}$$}

$$\eqalignno{\int f
&=\sup_{\Atop{g\text{ sederhana}}{g\le f\text{ hampir di mana-mana}}}\int g\cr
\displaycause{jika salah satunya berhingga, dan dengan demikian juga jika
salah satunya tak berhingga}
&=\sup_{\Atop{g\text{ sederhana}}{\Atop{g\le f\text{ hampir di mana-mana}}
{\mu F<\infty}}}\int_Fg
\le\sup_{\mu F<\infty}\int_Ff
\le\int f,\cr}$$

\noindent sehingga kita memperoleh kesamaan yang dicari.
}%end of proof of 213B

\leader{*213C}{Proposisi} Misalkan $(X,\Sigma,\mu)$ suatu ruang ukur
lengkap yang ditentukan secara lokal, dan misalkan $\mu^*$ ukuran luar yang
diturunkan dari $\mu$\cmmnt{ (132A-132B)}.   Maka ukuran yang diperoleh dari
$\mu^*$ dengan metode \Caratheodory adalah $\mu$ sendiri.

\proof{ Nyatakan $\check\mu$ sebagai ukuran yang diperoleh dengan metode
\Caratheodory dari $\mu^*$, dan domainnya sebagai $\check\Sigma$.

\medskip

{\bf (a)} Jika $E\in\Sigma$ dan $A\subseteq X$ maka
$\mu^*(A\cap E)+\mu^*(A\setminus E)=\mu^*A$ (132Af), jadi
$E\in\check\Sigma$.   Sekarang
$\check\mu E=\mu^*E=\mu E$ (132Ac).   Jadi $\Sigma\subseteq\check\Sigma$
dan $\mu=\check\mu\restr\Sigma$.

\medskip

{\bf (b)} Sekarang misalkan $H\in\check\Sigma$.   Ambil $E\in\Sigma$
dengan $\mu E<\infty$.   Maka $H\cap E\in\Sigma$.   \Prf\ Misalkan $E_1$,
$E_2\in\Sigma$ berturut-turut merupakan selubung terukur bagi $E\cap H$
dan $E\setminus H$, keduanya termuat dalam $E$ (132Ee).   Karena
$H\in\check\Sigma$,

\Centerline{$\mu E_1+\mu E_2=\mu^*(E\cap H)+\mu^*(E\setminus H)
=\mu^*E=\mu E$.}

\noindent Karena $E_1\cup E_2=E$,

\Centerline{$\mu(E_1\cap E_2)=\mu E_1+\mu E_2-\mu E=0$.}

\noindent Sekarang $E_1\setminus(E\cap H)\subseteq E_1\cap E_2$;  karena
$\mu$ lengkap, $E_1\setminus(E\cap H)$ dan $E\cap H$ termasuk dalam
$\Sigma$.\ \Qed

Karena $E$ sebarang dan $\mu$ ditentukan secara lokal, $H\in\Sigma$.
Karena $H$ sebarang, $\check\Sigma=\Sigma$ dan $\check\mu=\mu$.
}%end of proof of 213C

\leader{213D}{Versi c.l.d.: Proposisi} Misalkan $(X,\Sigma,\mu)$ suatu
ruang ukur.    Nyatakan $(X,\hat\Sigma,\hat\mu)$ sebagai
kelengkapannya\cmmnt{ (212C)} dan nyatakan $\Sigma^f$ sebagai
$\{E:E\in\Sigma,\,\mu E<\infty\}$.   Tetapkan

\Centerline{$\tilde\Sigma=\{H:H\subseteq X,\,H\cap E\in\hat\Sigma$ untuk
setiap $E\in\Sigma^f\}$,}

\noindent dan, untuk $H\in\tilde\Sigma$, tetapkan

\Centerline{$\tilde\mu H=\sup\{\hat\mu(H\cap E):E\in\Sigma^f\}$.}

\noindent Maka $(X,\tilde\Sigma,\tilde\mu)$ merupakan ruang ukur lengkap
yang ditentukan secara lokal.

\proof{{\bf (a)} Pertama-tama kita periksa bahwa $\tilde\Sigma$ merupakan
aljabar-sigma.   \Prf\ {(i)} $\emptyset\cap E=\emptyset\in\hat\Sigma$
untuk setiap $E\in\Sigma^f$, jadi $\emptyset\in\tilde\Sigma$.   {(ii)} Jika
$H\in\tilde\Sigma$, maka

\Centerline{$(X\setminus H)\cap E
=E\setminus(E\cap H)\in\hat\Sigma$}

\noindent untuk setiap $E\in\Sigma^f$, jadi $X\setminus H\in\tilde\Sigma$.
{(iii)} Jika $\sequencen{H_n}$ merupakan barisan dalam $\tilde\Sigma$
dengan gabungan $H$, maka

\Centerline{$H\cap E=\bigcup_{n\in\Bbb N}H_n\cap E\in\hat\Sigma$}

\noindent untuk setiap
$E\in\Sigma^f$, jadi $H\in\tilde\Sigma$.\ \Qed

\medskip

{\bf (b)} Jelas bahwa $\tilde\mu \emptyset=0$.   Jika
$\sequencen{H_n}$ merupakan barisan saling lepas dalam $\tilde\Sigma$
dengan gabungan $H$, maka

$$\eqalign{\tilde\mu H
&=\sup\{\hat\mu(H\cap E):E\in\Sigma^f\}\cr
&=\sup\{\sum_{n=0}^{\infty}\hat\mu(H_n\cap E):E\in\Sigma^f\}
\le\sum_{n=0}^{\infty}\tilde\mu H_n.\cr}$$

\noindent Sebaliknya, diberikan $a<\sum_{n=0}^{\infty}\tilde\mu H_n$,
terdapat $m\in\Bbb N$ sehingga $a<\sum_{n=0}^m\tilde\mu H_n$;  selanjutnya
dapat dipilih $E_0,\ldots,E_m\in\Sigma^f$ sedemikian sehingga
$a\le\sum_{n=0}^m\hat\mu(H_n\cap E_n)$.   Tetapkan $E=\bigcup_{n\le
m}E_n\in\Sigma^f$;  maka

\Centerline{$\tilde\mu H\ge\hat\mu(H\cap E)
=\sum_{n=0}^{\infty}\hat\mu(H_n\cap E)
\ge\sum_{n=0}^m\hat\mu(H_n\cap E_n)
\ge a$.}

\noindent Karena $a$ sebarang, $\tilde\mu
H\ge\sum_{n=0}^{\infty}\tilde\mu H_n$ dan $\tilde\mu
H=\sum_{n=0}^{\infty}\tilde\mu H_n$.

\medskip

{\bf (c)} Dengan demikian $(X,\tilde\Sigma,\tilde\mu)$ merupakan ruang
ukur.   Untuk membuktikan bahwa ruang ini semihingga, pertama perhatikan bahwa
$\hat\Sigma\subseteq\tilde\Sigma$ (karena jika $H\in\hat\Sigma$ maka
pasti $H\cap E\in\hat\Sigma$ untuk setiap $E\in\Sigma^f$), dan bahwa
$\tilde\mu H=\hat\mu H$ setiap kali $\hat\mu H<\infty$ (sebab menurut
definisi dalam 212Ca terdapat $E\in\Sigma^f$ sehingga
$H\subseteq E$, sehingga
$\tilde\mu H=\hat\mu(H\cap E)=\hat\mu H$).   Sekarang misalkan
$H\in\tilde\Sigma$ dan $\tilde\mu H=\infty$.   Tentu terdapat
$E\in\Sigma^f$ sehingga $\hat\mu(H\cap E)>0$;  maka $0<\hat\mu(H\cap
E)<\infty$, jadi $0<\tilde\mu(H\cap E)<\infty$.

\medskip

{\bf (d)} Jadi $(X,\tilde\Sigma,\tilde\mu)$ merupakan ruang ukur semihingga.
Untuk membuktikan bahwa ruang ini ditentukan secara lokal, misalkan
$H\subseteq X$ sedemikian sehingga
$H\cap G\in\tilde\Sigma$ setiap kali $G\in\tilde\Sigma$ dan
$\tilde\mu G<\infty$.   Secara khusus, harus berlaku
$H\cap E\in\tilde\Sigma$
untuk setiap $E\in\Sigma^f$.   Namun, ini justru berarti bahwa
$H\cap E\in\hat\Sigma$ untuk setiap $E\in\Sigma^f$, sehingga
$H\in\tilde\Sigma$.
Karena $H$ sebarang, $(X,\tilde\Sigma,\tilde\mu)$ ditentukan secara lokal.

\medskip

{\bf (e)} Untuk membuktikan bahwa $(X,\tilde\Sigma,\tilde\mu)$ lengkap,
misalkan $A\subseteq H\in\tilde\Sigma$ dan $\tilde\mu H=0$.
Untuk setiap $E\in\Sigma^f$ berlaku $\hat\mu(H\cap E)=0$.
Karena $(X,\hat\Sigma,\hat\mu)$ lengkap dan
$A\cap E\subseteq H\cap E$, maka $A\cap E\in\hat\Sigma$.   Karena $E$
sebarang, $A\in\tilde\Sigma$.
}%end of proof of 213D

\leader{213E}{Definisi} Untuk sebarang ruang ukur $(X,\Sigma,\mu)$,
$(X,\tilde\Sigma,\tilde\mu)$ yang dikonstruksi dalam 213D akan disebut
{\bf versi c.l.d.\ } (`versi lengkap yang ditentukan secara lokal') dari
$(X,\Sigma,\mu)$;  dan $\tilde\mu$ disebut {\bf versi c.l.d.\ } dari $\mu$.

\leader{213F}{}\cmmnt{ Dengan mengikuti pola yang sama seperti dalam
212E-212G, %212E 212F 212G
saya mulai dengan beberapa catatan elementer yang memudahkan penggunaan
konstruksi ini.

\medskip

\noindent}{\bf Proposisi} Misalkan $(X,\Sigma,\mu)$ suatu ruang ukur dan
$(X,\tilde\Sigma,\tilde\mu)$ versi c.l.d.\ yang bersesuaian.

(a) $\Sigma\subseteq\tilde\Sigma$ dan $\tilde\mu E=\mu E$ setiap kali
$E\in\Sigma$ dan $\mu E<\infty$ -- sebenarnya, jika $(X,\hat\Sigma,\hat\mu)$
adalah kelengkapan $(X,\Sigma,\mu)$, $\hat\Sigma\subseteq\tilde\Sigma$
dan $\tilde\mu E=\hat\mu E$ setiap kali $\hat\mu E<\infty$.

(b) Jika $\tilde\mu^*$ dan $\mu^*$ berturut-turut menyatakan ukuran luar
yang ditentukan oleh $\tilde\mu$ dan $\mu$, maka
$\tilde\mu^*A\le\mu^*A$ untuk setiap $A\subseteq X$, dengan kesamaan jika
$\mu^*A$ berhingga.   Secara khusus, setiap himpunan yang terabaikan terhadap
$\mu$ juga terabaikan terhadap $\tilde\mu$;  akibatnya, setiap himpunan yang
koterabaikan terhadap $\mu$ juga koterabaikan terhadap $\tilde\mu$.

(c) Jika $H\in\tilde\Sigma$,

\quad (i) $\tilde\mu H=\sup\{\mu F:F\in\Sigma$, $\mu F<\infty$,
$F\subseteq H\}$;

\quad (ii) terdapat $E\in\Sigma$ sehingga
$E\subseteq H$ dan $\mu E=\tilde\mu H$, sehingga
jika $\tilde\mu H<\infty$ maka $\tilde\mu(H\setminus E)=0$.

\proof{{\bf (a)} Pernyataan ini sudah tercakup dalam catatan pada bukti
213D.

\medskip

{\bf (b)} Jika $\mu^*A=\infty$ maka pasti $\tilde\mu^*A\le\mu^*A$.   Jika
$\mu^*A<\infty$, ambil $E\in\Sigma$ sedemikian rupa sehingga $A\subseteq E$ dan
$\mu E=\mu^*A$ (132Aa).   Kemudian

\Centerline{$\tilde\mu^*A\le\tilde\mu E=\mu E=\mu^*A$.}

\noindent Sebaliknya, jika $A\subseteq H\in\tilde\Sigma$, maka

\Centerline{$\tilde\mu H\ge\hat\mu(H\cap E)\ge\hat\mu^*A=\mu^*A$,}

\noindent dengan menggunakan 212Ea.   Jadi $\mu^*A\le\tilde\mu^*A$ dan
$\mu^*A=\tilde\mu^*A$.

\medskip

{\bf (c)} Nyatakan $\Sigma^f$ sebagai $\{E:E\in\Sigma,\,\mu E<\infty\}$;
maka menurut definisi dalam 213D,
$\tilde\mu H=\sup\{\hat\mu(H\cap E):E\in\Sigma^f\}$.   Ambil barisan
$\sequencen{E_n}$ dalam $\Sigma^f$ sedemikian sehingga
$\tilde\mu H=\sup_{n\in\Bbb N}\hat\mu(H\cap E_n)$.   Untuk setiap
$n\in\Bbb N$ ada $F_n\in\Sigma$ sehingga
$F_n\subseteq H\cap E_n$ dan
$\mu F_n=\hat\mu(H\cap E_n)$ (212C).   Tetapkan $E=\bigcup_{n\in\Bbb N}F_n$.
Maka $E\in\Sigma$, $E\subseteq H$ dan

\Centerline{$\tilde\mu H
=\sup_{n\in\Bbb N}\mu F_n
\le\lim_{n\to\infty}\mu(\bigcup_{i\le n}F_i)
=\mu E
=\lim_{n\to\infty}\tilde\mu(\bigcup_{i\le n}F_i)
\le\tilde\mu H,$}

\noindent jadi $\mu E=\tilde\mu H$, dan jika $\tilde\mu H<\infty$ maka
$\tilde\mu(H\setminus E)=0$.    Pada saat yang sama,

$$\eqalignno{\tilde\mu H
&=\sup_{n\in\Bbb N}\mu F_n
\le\sup_{F\in\Sigma^f,F\subseteq H}\mu F
=\sup_{F\in\Sigma^f,F\subseteq H}\tilde\mu F\cr
\displaycause{sekali lagi menurut (a)}
&\le\tilde\mu H,\cr}$$

\noindent sehingga di sini pun berlaku kesamaan.
}%end of proof of 213F

\leader{213G}{}\cmmnt{ Langkah berikutnya ialah menelaah fungsi-fungsi yang
terukur atau terintegralkan terhadap $\tilde\mu$.

\medskip

\noindent}{\bf Proposisi} Misalkan $(X,\Sigma,\mu)$ suatu ruang ukur dan
$(X,\tilde\Sigma,\tilde\mu)$ versi c.l.d.\ yang bersesuaian.

(a) Jika fungsi bernilai riil $f$ yang terdefinisi pada suatu subhimpunan
$X$ terukur secara virtual terhadap $\mu$, maka fungsi tersebut terukur terhadap
$\tilde\Sigma$.
%should we have $[-\infty,\infty]$-valued here?

(b) Jika suatu fungsi bernilai riil dapat diintegralkan terhadap $\mu$, maka ia
$\tilde\mu$-terintegralkan dengan integral yang sama.

(c) Jika $f$ merupakan fungsi bernilai riil yang dapat diintegralkan terhadap
$\tilde\mu$, maka terdapat fungsi bernilai riil yang dapat diintegralkan
terhadap $\mu$ dan sama dengan
$f\,\,\tilde\mu$-hampir di mana-mana.

\proof{ Nyatakan $\Sigma^f$ sebagai $\{E:E\in\Sigma,\,\mu E<\infty\}$.
Menurut 213Fa, $\tilde\mu$ dan $\mu$ berimpit pada $\Sigma^f$.

\medskip

{\bf (a)} Menurut 212Fa, $f$ terukur terhadap $\hat\Sigma$, dengan
$\hat\Sigma$ domain kelengkapan $\mu$;  karena
$\hat\Sigma\subseteq\tilde\Sigma$, $f$ adalah $\tilde\Sigma$-terukur.

\medskip

{\bf (b)(i)} Jika $f$ fungsi sederhana terhadap $\mu$, maka fungsi itu juga
fungsi sederhana terhadap $\tilde\mu$,
dan $\int fd\mu=\int fd\tilde\mu$, karena $\tilde\mu E=\mu E$ untuk setiap
$E\in\Sigma^f$.

\medskip

\quad{\bf (ii)} Jika $f$ nonnegatif dan dapat diintegralkan terhadap $\mu$,
terdapat barisan tak menurun $\sequencen{f_n}$ dari fungsi sederhana terhadap
$\mu$ yang konvergen ke $f\,\,\mu$-hampir di mana-mana;  menurut
213Fb, setiap himpunan yang terabaikan terhadap $\mu$ juga terabaikan terhadap
$\tilde\mu$, sehingga $\sequencen{f_n}$ konvergen ke
$f\,\,\tilde\mu$-hampir di mana-mana.   Berdasarkan teorema B.Levi (123A), $f$ dapat
diintegralkan terhadap $\tilde\mu$, dengan

\Centerline{$\int fd\tilde\mu=\lim_{n\to\infty}\int f_nd\tilde\mu
=\lim_{n\to\infty}\int f_nd\mu=\int fd\mu$.}

\medskip

\quad{\bf (iii)} Secara umum, jika $\int fd\mu$ didefinisikan dalam
$\Bbb R$, kita punya

\Centerline{$\int fd\tilde\mu
=\int f^+d\tilde\mu-\int f^-d\tilde\mu
=\int f^+d\mu-\int f^-d\mu=\int fd\mu$,}

\noindent dengan $f^+$ menyatakan $f\vee 0$ dan $f^-$ menyatakan
$(-f)\vee 0$.

\medskip

{\bf (c)(i)} Misalkan $f$ fungsi sederhana terhadap $\tilde\mu$.   Nyatakan
fungsi ini sebagai
$\sum_{i=0}^na_i\chi H_i$ dengan $\tilde\mu H_i<\infty$ untuk setiap $i$.
Pilih $E_0,\ldots,E_n\in\Sigma$ sehingga $E_i\subseteq H_i$ dan
$\tilde\mu(H_i\setminus E_i)=0$ untuk setiap $i$ (menggunakan 213Fc di atas).
Maka $g=\sum_{i=0}^na_i\chi E_i$ merupakan fungsi sederhana terhadap $\mu$,
$g=f\,\,\tilde\mu$-hampir di mana-mana, dan
$\int g\,d\mu=\int fd\tilde\mu$.

\medskip

\quad{\bf (ii)} Misalkan $f$ fungsi nonnegatif yang dapat diintegralkan
terhadap $\tilde\mu$.   Ambil barisan tak menurun $\sequencen{f_n}$ dari
fungsi-fungsi sederhana terhadap $\tilde\mu$ yang konvergen
$\tilde\mu$-hampir di mana-mana ke $f$.   Untuk setiap $n$, pilih fungsi
sederhana terhadap $\mu$, katakanlah $g_n$, yang sama $\tilde\mu$-hampir di mana-mana
dengan $f_n$.   Maka $\{x:g_{n+1}(x)<g_n(x)\}$ termasuk dalam $\Sigma^f$ dan
terabaikan terhadap $\tilde\mu$, sehingga juga terabaikan terhadap $\mu$.
Jadi $\sequencen{g_n}$ tak menurun hampir di mana-mana terhadap $\mu$.   Karena

\Centerline{$\lim_{n\to\infty}\int g_nd\mu
=\lim_{n\to\infty}\int f_nd\tilde\mu=\int fd\tilde\mu$,}

\noindent teorema B.Levi menunjukkan bahwa $\sequencen{g_n}$ konvergen
$\mu$-hampir di mana-mana ke suatu fungsi yang dapat diintegralkan terhadap
$\mu$, katakanlah $g$;  karena setiap himpunan yang terabaikan terhadap $\mu$ juga
terabaikan terhadap $\tilde\mu$,

$$\eqalign{(X\setminus\dom f)
&\cup(X\setminus\dom g)\cr
&\cup\bigcup_{n\in\Bbb N}\{x:f_n(x)\ne g_n(x)\}\cr
&\cup\{x:x\in\dom
f,\,f(x)\ne\sup_{n\in\Bbb N}f_n(x)\}\cr
&\cup\{x:x\in\dom
g,\,g(x)\ne\sup_{n\in\Bbb N}g_n(x)\}\cr}$$

\noindent terabaikan terhadap $\tilde\mu$, dan
$f=g\,\,\tilde\mu$-hampir di mana-mana.

\medskip

\quad{\bf (iii)} Jika $f$ dapat diintegralkan terhadap $\tilde\mu$, nyatakan
fungsi ini sebagai $f_1-f_2$, dengan $f_1$ dan $f_2$ nonnegatif serta dapat
diintegralkan terhadap $\tilde\mu$.   Terdapat fungsi yang dapat diintegralkan
terhadap $\mu$, yaitu $g_1$, $g_2$, sedemikian sehingga $f_1=g_1$ dan
$f_2=g_2\,\,\tilde\mu$-hampir di mana-mana;  jadi $g=g_1-g_2$ dapat
diintegralkan terhadap $\mu$ dan sama dengan $f\,\,\tilde\mu$-hampir di
mana-mana.
}%end of proof of 213G

\leader{213H}{}\cmmnt{ Ketiga, saya beralih ke pengaruh konstruksi ini
terhadap sifat-sifat lain yang dibahas dalam bab ini.

\medskip

\noindent}{\bf Proposisi} Misalkan $(X,\Sigma,\mu)$ suatu ruang ukur,
$(X,\hat\Sigma,\hat\mu)$ kelengkapannya, dan $(X,\tilde\Sigma,\tilde\mu)$
versi c.l.d.\ yang bersesuaian.

(a) Jika $(X,\Sigma,\mu)$ merupakan ruang probabilitas, berhingga total,
$\sigma$-hingga, atau dapat dilokalkan secara ketat, maka demikian pula
$(X,\tilde\Sigma,\tilde\mu)$, dan dalam semua kasus ini
$\tilde\mu=\hat\mu$;

(b) jika $(X,\Sigma,\mu)$ dapat dilokalkan, maka demikian pula
$(X,\tilde\Sigma,\tilde\mu)$, dan untuk setiap $H\in\tilde\Sigma$ ada
suatu $E\in\Sigma$ sehingga $\tilde\mu(E\symmdiff H)=0$;

(c) $(X,\Sigma,\mu)$ semihingga jika dan hanya jika
$\tilde\mu F=\mu F$ untuk setiap $F\in\Sigma$;  dalam hal ini
$\int fd\tilde\mu=\int fd\mu$ setiap kali integral yang terakhir terdefinisi
dalam $[-\infty,\infty]$;

(d) suatu himpunan $H\in\tilde\Sigma$ merupakan atom bagi $\tilde\mu$ jika
dan hanya jika terdapat atom $E$ bagi $\mu$ sedemikian sehingga
$\mu E<\infty$ dan
$\tilde\mu(H\symmdiff E)=0$;

(e) jika $(X,\Sigma,\mu)$ tanpa atom atau atomik murni, maka demikian pula
$(X,\tilde\Sigma,\tilde\mu)$;

(f) $(X,\Sigma,\mu)$ lengkap dan ditentukan secara lokal jika dan hanya jika
$\tilde\mu=\mu$.


\proof{{\bf (a)(i)} Pertama kita tunjukkan bahwa jika $(X,\Sigma,\mu)$
dapat dilokalkan secara ketat, maka
$\tilde\mu=\hat\mu$.   \Prf\ Misalkan $\langle X_i\rangle_{i\in I}$ suatu
dekomposisi $X$ bagi $\mu$;  maka keluarga ini juga merupakan dekomposisi
bagi $\hat\mu$ (212Gb).   Jika $H\in\tilde\Sigma$, berlaku
$H\cap X_i\in\hat\Sigma$ untuk setiap $i$, dan karenanya $H\in\hat\Sigma$;
lebih-lebih lagi,

$$\eqalign{\hat\mu H
&=\sum_{i\in I}\hat\mu(H\cap X_i)
=\sup\{\sum_{i\in J}\hat\mu(H\cap X_i):J\subseteq I
\text{ berhingga}\}\cr
&\le\sup\{\hat\mu(H\cap E):E\in\Sigma,\,\mu E<\infty\}
=\tilde\mu H
\le\hat\mu H.\cr}$$

\noindent Jadi $\hat\mu H=\tilde\mu H$ untuk setiap $H\in\tilde\Sigma$ dan
$\hat\mu=\tilde\mu$.\ \Qed

\medskip

\quad{\bf (ii)} Akibatnya, jika $(X,\Sigma,\mu)$ merupakan ruang
probabilitas, berhingga total, $\sigma$-hingga, atau dapat dilokalkan secara
ketat, maka demikian pula $(X,\tilde\Sigma,\tilde\mu)$;  gunakan 212Ga-212Gb
untuk melihat bahwa $(X,\hat\Sigma,\hat\mu)$ mempunyai sifat yang dimaksud.

\medskip

{\bf (b)} Jika $(X,\Sigma,\mu)$ dapat dilokalkan, ambil sebarang subhimpunan
$\Cal H$ dari $\tilde\Sigma$.   Tetapkan

\Centerline{$\Cal E=\{E:E\in\Sigma^f,\,\exists\enskip
H\in\Cal H,\,E\subseteq H\}$}

\noindent dengan $\Sigma^f=\{E:\mu E<\infty\}$ seperti biasa.
Di dalam $(X,\Sigma,\mu)$, pilih $F\in\Sigma$ sebagai supremum esensial
bagi $\Cal E$.

\medskip

\quad{\bf (i)} \Quer\ Andaikan, jika mungkin, terdapat
$H\in\Cal H$ sehingga $\tilde\mu(H\setminus F)>0$.   Oleh 213F(c-i),
ada
$E\in\Sigma^f$ sehingga $E\subseteq H\setminus F$ dan $\mu E>0$.
$E$ ini termasuk dalam $\Cal E$ dan $\mu(E\setminus F)=\mu E>0$;  hal ini
mustahil karena $F$ merupakan supremum esensial bagi $\Cal E$.\ \Bang

\medskip

\quad{\bf (ii)} Jadi $\tilde\mu(H\setminus F)=0$ untuk setiap $H\in\Cal H$.
Sekarang ambil sebarang $G\in\tilde\Sigma$ sedemikian sehingga
$\tilde\mu(H\setminus G)=0$ untuk setiap $H\in\Cal H$.   \Quer\ Jika
$\tilde\mu(F\setminus G)>0$, terdapat $E_0\in\Sigma^f$ sehingga
$E_0\subseteq F\setminus G$ dan
$\mu E_0>0$.   Jika $E\in\Cal E$, terdapat
$H\in\Cal H$ sehingga $E\subseteq H$, sehingga
$E\setminus(F\setminus E_0)\subseteq H\setminus(F\cap G)$, sedangkan
$\mu(E\setminus(F\setminus E_0))<\infty$;  jadi

\Centerline{$\mu(E\setminus(F\setminus E_0))
\le\tilde\mu(H\setminus(F\cap G))
\le\tilde\mu(H\setminus F)+\tilde\mu(H\setminus G)
=0$.}

\noindent Karena $F$ merupakan supremum esensial bagi $\Cal E$ dalam $\Sigma$,

\Centerline{$0
=\mu(F\setminus(F\setminus E_0))
=\mu E_0$.  \Bang}

\noindent Ini menunjukkan bahwa $F$ merupakan supremum esensial bagi $\Cal H$
dalam $\tilde\Sigma$.   Karena $\Cal H$ sebarang,
$(X,\tilde\Sigma,\tilde\mu)$ dapat dilokalkan.

\medskip

\quad{\bf (iii)} Argumen dalam (i)-(ii) sebenarnya menunjukkan bahwa jika
$\Cal H\subseteq\tilde\Sigma$, maka $\Cal H$ memiliki supremum esensial $F$
dalam $\tilde\Sigma$ sedemikian sehingga $F$ bahkan termasuk dalam $\Sigma$.
Dengan mengambil
$\Cal H=\{H\}$, kita melihat bahwa jika $H\in\tilde\Sigma$ terdapat
$F\in\Sigma$ sehingga $\tilde\mu(H\symmdiff F)=0$.

\medskip

{\bf (c)} Dari 213Fa kita telah mengetahui bahwa
$\tilde\mu E\le\mu E$ untuk setiap $E\in\Sigma$, dengan kesamaan jika
$\mu E<\infty$.

\medskip

\quad{\bf (i)} Jika $(X,\Sigma,\mu)$ semihingga, maka 213A dan 213F(c-i)
menunjukkan bahwa untuk sebarang $F\in\Sigma$ berlaku

\Centerline{$\mu F
=\sup\{\mu E:E\in\Sigma,\,E\subseteq F,\,\mu E<\infty\}
=\tilde\mu F$.}

\medskip

\quad{\bf (ii)} Misalkan $\tilde\mu F=\mu F$ untuk setiap $F\in\Sigma$.
Jika $\mu F=\infty$, maka $\tilde\mu F=\infty$, sehingga (sekali lagi oleh
213F(c-i)) terdapat $E\in\Sigma^f$ dengan $E\subseteq F$ dan $\mu E>0$;
karena $F$ sebarang, $(X,\Sigma,\mu)$ semihingga.

\medskip

\quad{\bf (iii)} Jika $f$ nonnegatif dan $\int fd\mu=\infty$, maka $f$
terukur secara virtual terhadap $\mu$, sehingga terukur terhadap
$\tilde\Sigma$ (213Ga), dan terdefinisi hampir di mana-mana terhadap $\mu$,
sehingga juga terhadap $\tilde\mu$.   Sekarang

$$\eqalign{\int fd\tilde\mu
&=\sup\{\int g\,d\tilde\mu:g\text{ merupakan fungsi }\tilde\mu\text{-sederhana},\,
   0\le g\le f\,\,\tilde\mu\text{-hampir di mana-mana}\}\cr
&\ge\sup\{\int g\,d\mu:g\text{ merupakan fungsi }\mu\text{-sederhana},\,
   0\le g\le f\,\,\mu\text{-hampir di mana-mana}\}=\infty\cr}$$

\noindent menurut 213B.   Bersama 213Gb, ini menunjukkan bahwa
$\int fd\tilde\mu=\int fd\mu$ setiap kali $f$ nonnegatif dan
$\int fd\mu$ terdefinisi dalam $[0,\infty]$.   Dengan menerapkannya pada
bagian positif dan negatif $f$, kita memperoleh
$\int fd\tilde\mu=\int fd\mu$ setiap kali integral yang terakhir terdefinisi
dalam $[-\infty,\infty]$.

\medskip

{\bf (d)(i)} Jika $H\in\tilde\Sigma$ merupakan atom bagi $\tilde\mu$, maka
ada $E\in\Sigma^f$ sehingga $E\subseteq H$ dan $0<\mu E<\infty$.
Dalam hal ini, $\tilde\mu E>0$, sehingga $\tilde\mu(H\setminus E)$ harus nol.
Jika $F\in\Sigma$ dan $F\subseteq E$, maka
$\mu F=\tilde\mu F=0$ atau $\mu(E\setminus F)=\tilde\mu(H\setminus F)=0$.
Jadi $E\in\Sigma$ adalah atom untuk $\mu$ dengan
$\tilde\mu(H\symmdiff E)=0$ dan $\mu E<\infty$.

\medskip

\quad{\bf (ii)} Jika $H\in\tilde\Sigma$ dan terdapat atom $E$ bagi $\mu$
sedemikian sehingga $\mu E<\infty$ dan $\tilde\mu(H\symmdiff E)=0$, ambil
$G\in\tilde\Sigma$ sebagai subhimpunan $H$ dengan $\tilde\mu G>0$.
Kita mempunyai $\tilde\mu(E\cap G)=\tilde\mu(H\cap G)>0$, sehingga terdapat
$F\in\Sigma$ sehingga $F\subseteq E\cap G$ dan $\mu F>0$.   Sekarang
$\mu(E\setminus F)$ harus nol, jadi

\Centerline{$\tilde\mu(H\setminus G)\le\tilde\mu(H\setminus F)
=\tilde\mu(E\setminus F)=\mu(E\setminus F)=0$.}

\noindent Karena $G$ sebarang, $H$ merupakan atom bagi $\tilde\mu$.

\medskip

{\bf (e)} Jika $(X,\Sigma,\mu)$ tanpa atom, maka
$(X,\tilde\Sigma,\tilde\mu)$ juga tanpa atom, menurut (d).

Jika $(X,\Sigma,\mu)$ atomik murni, $H\in\tilde\Sigma$ dan
$\tilde\mu H>0$, maka terdapat $E\in\Sigma^f$ sehingga
$E\subseteq H$ dan $\mu E>0$.
Terdapat atom $F$ bagi $\mu$ dengan $F\subseteq E$;  karena $\mu F<\infty$,
menurut (d), $F$ merupakan atom bagi $\tilde\mu$.   Selain itu $F\subseteq H$.
Karena $H$ sebarang, $(X,\tilde\Sigma,\tilde\mu)$ atomik murni.

\medskip

{\bf (f)} Jika $\mu=\tilde\mu$, tentu $(X,\Sigma,\mu)$ harus
lengkap dan ditentukan secara lokal, karena $(X,\tilde\Sigma,\tilde\mu)$
lengkap dan ditentukan secara lokal.   Sebaliknya, jika $(X,\Sigma,\mu)$
lengkap dan ditentukan secara lokal, maka $\hat\mu=\mu$, sehingga (dengan
definisi dalam 213D) $\tilde\Sigma\subseteq\Sigma$ dan $\tilde\mu=\mu$,
menurut (c) di atas.
}%end of proof of 213H

\leader{213I}{Himpunan terabaikan yang ditentukan secara
lokal\dvrocolon{}}\cmmnt{ Gagasan sederhana berikut kadang-kadang berguna.

\medskip

\noindent}{\bf Definisi} Suatu ruang ukur $(X,\Sigma,\mu)$ dikatakan mempunyai
{\bf himpunan terabaikan yang ditentukan secara lokal} jika untuk setiap
$A\subseteq X$ yang tidak terabaikan terdapat $E\in\Sigma$ dengan
$\mu E<\infty$ sedemikian sehingga $A\cap E$ tidak terabaikan.

\leader{213J}{Proposisi} Jika suatu ruang ukur $(X,\Sigma,\mu)$ {\it entah}
dapat dilokalkan secara ketat {\it atau} lengkap dan ditentukan secara lokal,
maka ruang tersebut mempunyai himpunan terabaikan yang ditentukan secara lokal.

\proof{ Misalkan $A\subseteq X$ sedemikian sehingga $A\cap E$ terabaikan
setiap kali $\mu E<\infty$;  kita harus menunjukkan bahwa $A$ terabaikan.

\medskip

{\bf (i)} Jika $\mu$ dapat dilokalkan secara ketat, misalkan
$\familyiI{X_i}$ suatu dekomposisi $X$.   Untuk setiap $i\in I$,
$A\cap X_i$ terabaikan, sehingga dapat dipilih himpunan terabaikan
$E_i\in\Sigma$ dengan $A\cap X_i\subseteq E_i$.   Tetapkan
$E=\bigcup_{i\in I}E_i\cap X_i$.   Maka
$\mu E=\sum_{i\in I}\mu(E_i\cap X_i)=0$ dan $A\subseteq E$, jadi $A$
terabaikan.

\medskip

{\bf (ii)} Jika $\mu$ lengkap dan ditentukan secara lokal, ambil sebarang
himpunan terukur $E$ berukuran berhingga.   Maka $A\cap E$ terabaikan dan
karena itu terukur;  karena $E$ sebarang, $A$ terukur;  karena $\mu$
semihingga, $A$ terabaikan.
}%end of proof of 213J

\leader{*213K}{Lema} Jika suatu ruang ukur $(X,\Sigma,\mu)$ mempunyai
himpunan terabaikan yang ditentukan secara lokal, dan
$\Cal E\subseteq\Sigma$ mempunyai supremum esensial
$H\in\Sigma$\cmmnt{ dalam arti 211G}, maka
$H\setminus\bigcup\Cal E$ terabaikan.

\proof{ Tetapkan $A=H\setminus\bigcup\Cal E$.   Ambil sebarang
$F\in\Sigma$ dengan $\mu F<\infty$.   Maka $F\cap A$ mempunyai suatu
selubung terukur, katakanlah $V$ (sekali lagi menurut 132Ee).
Jika $E\in\Cal E$, maka

\Centerline{$\mu(E\setminus(X\setminus V))
=\mu(E\cap V)=\mu^*(E\cap F\cap A)=0$,}

\noindent sehingga $H\cap V=H\setminus(X\setminus V)$ terabaikan dan
$F\cap A$ terabaikan.   Karena $F$ sebarang dan $\mu$ mempunyai himpunan
terabaikan yang ditentukan secara lokal, $A$ terabaikan, sebagaimana dinyatakan.
}%end of proof of 213K

\leader{213L}{Proposisi} Misalkan $(X,\Sigma,\mu)$ suatu ruang ukur yang
dapat dilokalkan dan mempunyai himpunan terabaikan yang ditentukan secara
lokal. Maka setiap subhimpunan $A$ dari $X$ mempunyai selubung terukur.

\proof{ Tetapkan

\Centerline{$\Cal E=\{E:E\in\Sigma,\,\mu^*(A\cap E)=\mu E<\infty\}$.}

\noindent Misalkan $G$ suatu supremum esensial bagi $\Cal E$ dalam $\Sigma$.

\medskip

\quad{\bf (i)} $A\setminus G$ terabaikan.   \Prf\ Ambil sebarang himpunan
$F$ berukuran berhingga terhadap $\mu$.   Misalkan $E$ suatu selubung terukur
bagi $A\cap F$.   Maka $E\in\Cal E$, sehingga $E\setminus G$ terabaikan.
Karena $F\cap A\setminus G\subseteq E\setminus G$, maka
$F\cap A\setminus G$ terabaikan.   Karena $\mu$ mempunyai himpunan terabaikan
yang ditentukan secara lokal, ini cukup untuk menunjukkan bahwa
$A\setminus G$ terabaikan.\ \Qed

\medskip

\quad{\bf (ii)} Ambil himpunan terukur terabaikan $E_0$ yang memuat
$A\setminus G$, dan tetapkan $\tilde G=E_0\cup G$, sehingga
$\tilde G\in\Sigma$,
$A\subseteq\tilde G$ dan $\mu(\tilde G\setminus G)=0$.   \Quer\
Andaikan, jika mungkin, terdapat $F\in\Sigma$ sedemikian sehingga
$\mu^*(A\cap F)<\mu(\tilde G\cap F)$.   Misalkan $F_1\subseteq F$ suatu
selubung terukur bagi $A\cap F$.   Tetapkan
$H=X\setminus(F\setminus F_1)$;  maka $A\subseteq H$.   Jika
$E\in\Cal E$, maka

\Centerline{$\mu E=\mu^*(A\cap E)\le\mu(H\cap E)$,}

\noindent sehingga $E\setminus H$ terabaikan;  karena $E$ sebarang,
$G\setminus H$ dan $\tilde G\setminus H$ terabaikan.
Namun $\tilde G\cap F\setminus F_1\subseteq\tilde G\setminus H$ dan

\Centerline{$\mu(\tilde G\cap F\setminus F_1)
=\mu(\tilde G\cap F)-\mu^*(A\cap F)>0$. \Bang}

Ini menunjukkan bahwa $\tilde G$ merupakan selubung terukur bagi $A$,
sebagaimana diperlukan.
}%end of proof of 213L

\leader{213M}{Akibat} (a) Jika $(X,\Sigma,\mu)$ $\sigma$-hingga, maka
setiap subhimpunan dari $X$ mempunyai selubung terukur terhadap $\mu$.

(b) Jika $(X,\Sigma,\mu)$ dapat dilokalkan, maka setiap subhimpunan dari $X$
mempunyai selubung terukur terhadap versi c.l.d.\ dari $\mu$.

\proof{{\bf (a)} Gunakan 132Ee, atau 213L, 211Lc dan 213J.

\medskip

{\bf (b)} Gunakan 213L dan fakta bahwa versi c.l.d.\ dari $\mu$ dapat
dilokalkan, selain lengkap dan ditentukan secara lokal (213Hb).
}%end of proof of 213M

\leader{213N}{}\cmmnt{ Ketika kelak kita menggunakan konsep
`keterlokalan', konsep itu sering muncul melalui sifat berikut, yang
sesungguhnya mencirikan ruang yang dapat dilokalkan (213Xm).

\medskip

\noindent}{\bf Teorema} Misalkan $(X,\Sigma,\mu)$ suatu ruang ukur yang
dapat dilokalkan.   Misalkan $\Phi$ keluarga fungsi terukur bernilai riil,
yang semuanya terdefinisi pada subhimpunan terukur dari $X$, sedemikian
sehingga untuk setiap $f$, $g\in\Phi$, berlaku $f=g$ hampir di mana-mana pada
$\dom f\cap\dom g$. Maka terdapat fungsi terukur $h:X\to\Bbb R$ sedemikian
sehingga setiap $f\in\Phi$ sama dengan $h$ hampir di mana-mana pada $\dom f$.

\proof{ Untuk $q\in\Bbb Q$, $f\in\Phi$, tetapkan

\Centerline{$E_{fq}=\{x:x\in\dom f,\,f(x)\ge q\}\in\Sigma$.}

\noindent Untuk setiap $q\in\Bbb Q$, misalkan $E_q$ suatu supremum esensial
bagi $\{E_{fq}:f\in\Phi\}$ dalam $\Sigma$.   Tetapkan

\Centerline{$h^*(x)=\sup\{q:q\in\Bbb Q,\,x\in E_q\}\in[-\infty,\infty]$}

\noindent untuk $x\in X$, dengan mengambil $\sup\emptyset=-\infty$ bila perlu.

Jika $f$, $g\in\Phi$ dan $q\in\Bbb Q$, maka

$$\eqalign{E_{fq}\setminus(X\setminus(\dom g\setminus E_{gq}))
&=E_{fq}\cap\dom g\setminus E_{gq}\cr
&\subseteq\{x:x\in\dom f\cap\dom g,\,f(x)\ne g(x)\}\cr}$$

\noindent terabaikan;  karena $f$ sebarang,

\Centerline{$E_q\cap\dom g\setminus E_{gq}=E_q\setminus(X\setminus(\dom
g\setminus E_{gq}))$}

\noindent terabaikan.   Selain itu $E_{gq}\setminus E_q$ terabaikan, sehingga
$E_{gq}\symmdiff(E_q\cap\dom g)$ terabaikan.   Tetapkan
$H_g=\bigcup_{q\in\Bbb Q}E_{gq}\symmdiff(E_q\cap\dom g)$;  maka $H_g$
terabaikan.   Jika $x\in\dom g\setminus H_g$, maka untuk setiap
$q\in\Bbb Q$, $x\in E_q\iff x\in E_{gq}$;  akibatnya, untuk $x$ demikian,
$h^*(x)=g(x)$.   Jadi $h^*=g$ hampir di mana-mana pada $\dom g$;  dan ini
berlaku untuk setiap $g\in\Phi$.

Fungsi $h^*$ belum tentu bernilai riil.   Namun fungsi ini terukur, sebab
\Centerline{$\{x:h^*(x)>a\}=\bigcup\{E_q:q\in\Bbb Q,\,q>a\}\in\Sigma$}

\noindent untuk setiap bilangan riil $a$.   Jadi, jika kita memodifikasinya
dengan menetapkan

$$\eqalign{h(x)&=h^*(x)\text{ jika }h^*(x)\in\Bbb R,\cr
&=0\text{ jika }h^*(x)\in\{-\infty,\infty\},\cr}$$

\noindent kita memperoleh fungsi terukur bernilai riil $h:X\to\Bbb R$.
Untuk sebarang $g\in\Phi$, $h(x)$ sama dengan $g(x)$ setidaknya setiap kali
$h^*(x)=g(x)$, dan hal ini berlaku untuk hampir setiap $x\in\dom g$.   Jadi
$h$ merupakan fungsi yang diminta.
}%end of proof of 213N

\leader{213O}{}\cmmnt{ Berikut adalah suatu kriteria menarik dan berguna
agar sebuah ruang dapat dilokalkan secara ketat;  saya memperkenalkannya di
sini, walaupun kriteria ini jarang digunakan dalam jilid ini.

\medskip

\noindent}{\bf Proposisi} Misalkan $(X,\Sigma,\mu)$ suatu ruang lengkap yang
ditentukan secara lokal.

(a) Misalkan terdapat keluarga saling lepas $\Cal E\subseteq\Sigma$
sedemikian rupa sehingga ($\alpha$) $\mu E<\infty$ untuk setiap $E\in\Cal E$
($\beta$) setiap kali $F\in\Sigma$ dan $\mu F>0$, terdapat $E\in\Cal E$
sedemikian sehingga $\mu(E\cap F)>0$.   Maka $(X,\Sigma,\mu)$ dapat
dilokalkan secara ketat, $\bigcup\Cal E$ koterabaikan, dan
$\Cal E\cup\{X\setminus\bigcup\Cal E\}$ adalah
dekomposisi $X$.

(b) Misalkan $\langle X_i\rangle_{i\in I}$ suatu partisi $X$ menjadi
himpunan-himpunan terukur berukuran berhingga sedemikian sehingga setiap kali
$E\in\Sigma$ dan
$\mu E>0$ terdapat $i\in I$ sehingga $\mu(E\cap X_i)>0$.   Lalu
$(X,\Sigma,\mu)$ dapat dilokalkan secara ketat, dan
$\langle X_i\rangle_{i\in I}$ merupakan dekomposisi $X$.

\proof{{\bf (a)(i)}
Pertama, jika $F\in\Sigma$ dan $\mu F<\infty$, terdapat subkeluarga
terhitung $\Cal E'\subseteq\Cal E$ sedemikian sehingga
$\mu(F\setminus\bigcup\Cal E')=0$.   \Prf\ Tetapkan

\Centerline{$\Cal E'_n=\{E:E\in\Cal E,\,\mu(F\cap E)\ge 2^{-n}\}$ untuk
setiap $n\in\Bbb N$,}

\Centerline{$\Cal E'=\bigcup_{n\in\Bbb N}\Cal E'_n=\{E:E\in\Cal
E,\,\mu(F\cap E)>0\}$.}

\noindent Karena $\Cal E$ saling lepas, harus berlaku

\Centerline{$\#(\Cal E'_n)\le 2^n\mu F$}

\noindent untuk setiap $n\in\Bbb N$, sehingga setiap $\Cal E'_n$ berhingga,
dan $\Cal E'$, sebagai gabungan suatu barisan himpunan terhitung, juga
terhitung.   Tetapkan $E'=\bigcup\Cal E'$ dan $F'=F\setminus E'$, sehingga
$E'$ dan $F'$ termasuk dalam $\Sigma$.   Jika $E\in\Cal E'$, maka
$E\subseteq E'$, sehingga $\mu(E\cap F')=\mu\emptyset=0$;  jika $E\in\Cal
E\setminus\Cal E'$, maka $\mu(E\cap F')=\mu(E\cap F)=0$.   Jadi
$\mu(E\cap F')=0$ untuk setiap $E\in\Cal E$.   Berdasarkan hipotesis
($\beta$) pada $\Cal E$, $\mu F'=0$, sehingga
$\mu(F\setminus\bigcup\Cal E')=0$, sebagaimana diperlukan.\ \Qed

\medskip

\quad{\bf (ii)} Sekarang misalkan $H\subseteq X$ sedemikian sehingga
$H\cap E\in\Sigma$ untuk setiap $E\in\Cal E$.   Dalam hal ini
$H\in\Sigma$.   \Prf\ Ambil $F\in\Sigma$ dengan $\mu F<\infty$.   Misalkan
$\Cal E'\subseteq\Cal E$ suatu subkeluarga terhitung sedemikian sehingga
$\mu(F\setminus E')=0$, dengan $E'=\bigcup\Cal E'$.  Kemudian
$H\cap(F\setminus E')\in\Sigma$, karena $(X,\Sigma,\mu)$ lengkap.   Selain itu,
$H\cap
E'=\bigcup_{E\in\Cal E'}H\cap E\in\Sigma$.   Jadi

\Centerline{$H\cap F=(H\cap(F\setminus E'))\cup(F\cap(H\cap
E'))\in\Sigma$.}

\noindent Karena $F$ sebarang dan $(X,\Sigma,\mu)$ ditentukan secara lokal,
$H\in\Sigma$.\ \Qed

\medskip

\quad{\bf (iii)} Kita juga memperoleh bahwa
$\mu H=\sum_{E\in\Cal E}\mu(H\cap E)$ untuk setiap $H\in\Sigma$.
\Prf\ {($\alpha$)} Karena $\Cal E$ saling lepas,
$\sum_{E\in\Cal E'}\mu(H\cap E)\le\mu H$ untuk setiap subkeluarga berhingga
$\Cal E'\subseteq\Cal E$, jadi

\Centerline{$\sum_{E\in\Cal E}\mu(H\cap E)
=\sup\{\sum_{E\in\Cal E'}
  \mu(H\cap E):\Cal E'\subseteq\Cal E$ berhingga$\}\le\mu H$.}

\noindent ($\beta$) Untuk ketaksamaan sebaliknya, pertama tinjau kasus
$\mu H<\infty$.   Dengan (i), terdapat subkeluarga terhitung $\Cal E'\subseteq\Cal E$
sedemikian rupa sehingga $\mu(H\setminus\bigcup\Cal E')=0$, sehingga

\Centerline{$\mu H=\mu(H\cap\bigcup\Cal E')
=\sum_{E\in\Cal E'}\mu(H\cap E)\le\sum_{E\in\Cal E}\mu(H\cap E)$.}

\noindent{($\gamma$)} Secara umum, karena $(X,\Sigma,\mu)$ semihingga,

$$\eqalign{\mu H
&=\sup\{\mu F:F\subseteq H,\,\mu F<\infty\}\cr
&\le\sup\{\sum_{E\in\Cal E}\mu(F\cap E):F\subseteq H,\,\mu F<\infty\}
\le\sum_{E\in\Cal E}\mu(H\cap E).\cr}$$

\noindent Jadi, dalam semua kasus,
$\mu H\le\sum_{E\in\Cal E}\mu(H\cap E)$, dan kedua ruas sama.\ \Qed

\medskip

\quad{\bf (iv)} Secara khusus, dengan menetapkan $E_0=X\setminus\bigcup\Cal E$,
$E_0\in\Sigma$ dan $\mu E_0=0$;  yaitu, $\bigcup\Cal E$ adalah
koterabaikan.
Pertimbangkan $\Cal E^*=\Cal E\cup\{E_0\}$.
Keluarga ini merupakan partisi $X$ menjadi himpunan berukuran berhingga (kini menggunakan
hipotesis ($\alpha$) pada $\Cal E$).   Jika $H\subseteq X$ sedemikian rupa sehingga
$H\cap E\in\Sigma$ untuk setiap $E\in\Cal E^*$, maka $H\in\Sigma$ dan

\Centerline{$\mu H=\sum_{E\in\Cal E}\mu(H\cap E)
=\sum_{E\in\Cal E^*}\mu(H\cap E)$.}

\noindent Jadi $\Cal E^*$ (atau, jika dikehendaki, keluarga berindeks
$\langle E\rangle_{E\in\Cal E^*}$) merupakan dekomposisi yang membuktikan
bahwa $(X,\Sigma,\mu)$ dapat dilokalkan secara ketat.

\medskip

{\bf (b)} Terapkan (a) dengan $\Cal E=\{X_i:i\in I\}$, sambil memperhatikan
bahwa $E_0$ dalam (iv) kosong sehingga dapat dihilangkan.
}%end of proof of 213O

\exercises{
\leader{213X}{Latihan dasar (a)}
%\spheader 213Xa
Misalkan $(X,\Sigma,\mu)$ sebarang ruang ukur, $\mu^*$ ukuran luar yang
ditentukan oleh $\mu$, dan $\check\mu$ ukuran yang ditentukan dengan metode
\Caratheodory dari $\mu^*$;  nyatakan $\check\Sigma$ sebagai domain
$\check\mu$.   Tunjukkan bahwa
(i) $\check\mu$ memperluas kelengkapan $\hat\mu$ dari $\mu$;
(ii) jika $H\subseteq X$ sedemikian rupa sehingga
$H\cap F\in\check\Sigma$ setiap kali $F\in\Sigma$ dan $\mu F<\infty$, maka
$H\in\check\Sigma$;
(iii) $(\check\mu)^*=\mu^*$, sehingga fungsi yang dapat diintegralkan terhadap
$\check\mu$ dan $\mu$ sama (212Xb);  (iv) jika $\mu$ dapat dilokalkan secara
ketat, maka $\check\mu=\hat\mu$;  (v) jika $\mu$ ditentukan dengan metode
\Caratheodory dari suatu ukuran luar lain, maka $\mu=\check\mu$
(lih.\ 132Xa).
%213C

\sqheader 213Xb
Misalkan $\mu$ ukuran cacah yang dibatasi pada aljabar-sigma
terhitung-koterhitung dari suatu himpunan $X$ (211R, 211Ye).
(i) Tunjukkan bahwa versi c.l.d.\ $\tilde\mu$ dari $\mu$ tidak lain adalah
ukuran cacah pada $X$.   (ii) Tunjukkan bahwa $\check\mu$, sebagaimana
didefinisikan dalam 213Xa, sama dengan $\tilde\mu$, dan khususnya memperluas
kelengkapan $\mu$ secara ketat jika $X$ tak terhitung.
%213E

\spheader 213Xc Misalkan $(X,\Sigma,\mu)$ sebarang ruang ukur.   Untuk
$E\in\Sigma$, tetapkan

\Centerline{$\mu_{\text{sf}}E
=\sup\{\mu(E\cap F):F\in\Sigma,\,\mu F<\infty\}$.}

\quad{(i)} Tunjukkan bahwa $(X,\Sigma,\mu_{\text{sf}})$ merupakan ruang ukur
semihingga, dan sama dengan $(X,\Sigma,\mu)$ jika dan hanya jika
$(X,\Sigma,\mu)$ semihingga.

\quad{(ii)} Tunjukkan bahwa fungsi bernilai riil yang dapat diintegralkan
terhadap $\mu$, katakanlah $f$, juga dapat diintegralkan terhadap
$\mu_{\text{sf}}$, dengan
integral yang sama.

\quad(iii) Tunjukkan bahwa jika $E\in\Sigma$ maka $E$
dapat dinyatakan sebagai $E_1\cup E_2$ dengan $E_1$, $E_2\in\Sigma$,
$\mu E_1=\mu_{\text{sf}}E_1$ dan $\mu_{\text{sf}}E_2=0$.

\quad{(iv)} Tunjukkan bahwa jika $f$ fungsi bernilai riil yang dapat
diintegralkan terhadap $\mu_{\text{sf}}$ pada $X$, maka fungsi ini sama
hampir di mana-mana terhadap $\mu_{\text{sf}}$ dengan suatu fungsi yang dapat diintegralkan
terhadap $\mu$.

\quad{(v)} Tunjukkan bahwa jika $(X,\Sigma,\mu_{\text{sf}})$ lengkap, maka
$(X,\Sigma,\mu)$ juga lengkap.

\quad{(vi)} Tunjukkan bahwa $\mu$ dan $\mu_{\text{sf}}$ mempunyai versi
c.l.d.\ yang sama.
%213F

\spheader 213Xd Misalkan $(X,\Sigma,\mu)$ sebarang ruang ukur.   Definisikan
$\check\mu$ seperti dalam 213Xa.   Tunjukkan bahwa
$(\check\mu)_{\text{sf}}$, sebagaimana dikonstruksi dalam 213Xc, tepat
merupakan versi c.l.d.\ $\tilde\mu$ dari $\mu$, sehingga
$\check\mu=\tilde\mu$ jika dan hanya jika $\check\mu$ semihingga.
%213F, 213Xa, 213Xc

\spheader 213Xe Misalkan $(X,\Sigma,\mu)$ suatu ruang ukur.   Untuk
$A\subseteq X$, tetapkan
$\mu_* A=\sup\{\mu E:E\in\Sigma,\,\mu E<\infty,\,E\subseteq A\}$, seperti pada
113Yh.   (i) Tunjukkan bahwa ukuran yang dikonstruksi dari $\mu_*$ dengan
metode 113Yg/212Ya tidak lain adalah versi c.l.d.\ $\tilde\mu$ dari $\mu$.
(ii) Tunjukkan bahwa $\tilde\mu_*=\mu_*$.   (iii) Tunjukkan bahwa jika $\nu$
suatu ukuran lain pada $X$, dengan domain $\Tau$, maka
$\tilde\mu=\tilde\nu$ jika dan hanya jika
$\mu_*=\nu_*$.
%213F

\spheader 213Xf Misalkan $X$ suatu himpunan dan $\theta$ suatu ukuran luar
pada $X$. Tunjukkan bahwa $\theta_{\text{sf}}$, yang didefinisikan dengan

\Centerline{$\theta_{\text{sf}}A=\sup\{\theta B:B\subseteq A,\,\theta
B<\infty\}$}

\noindent juga merupakan ukuran luar pada $X$.   Tunjukkan bahwa ukuran yang
ditentukan dengan metode \Caratheodory dari $\theta$ dan
$\theta_{\text{sf}}$ mempunyai domain yang sama.
%213F

\spheader 213Xg Misalkan $(X,\Sigma,\mu)$ sebarang ruang ukur.   Tetapkan

\Centerline{$\mu^*_{\text{sf}}A
=\sup\{\mu^*(A\cap E):E\in\Sigma,\,\mu E<\infty\}$}

\noindent untuk setiap $A\subseteq X$.

\quad{(i)} Tunjukkan bahwa

\Centerline{$\mu^*_{\text{sf}}A=\sup\{\mu^*B:B\subseteq A,\,\mu^*B<\infty\}$}

\noindent untuk setiap $A$.

\quad{(ii)} Tunjukkan bahwa $\mu^*_{\text{sf}}$ adalah ukuran luar.

\quad(iii) Tunjukkan bahwa jika $A\subseteq X$ dan
$\mu^*_{\text{sf}}A<\infty$, terdapat $E\in\Sigma$ sehingga
$\mu^*_{\text{sf}}A=\mu^*(A\cap E)=\mu E$ dan
$\mu^*_{\text{sf}}(A\setminus E)=0$.   \Hint{ambil barisan tak menurun
$\sequencen{E_n}$ dari himpunan terukur berukuran berhingga sedemikian sehingga
$\mu^*_{\text{sf}}A=\lim_{n\to\infty}\mu^*(A\cap E_n)$, lalu ambil
$E\subseteq\bigcup_{n\in\Bbb N}E_n$ sebagai selubung terukur bagi
$A\cap\bigcup_{n\in\Bbb N}E_n$.}

\quad{(iv)} Tunjukkan bahwa ukuran yang ditentukan dari
$\mu^*_{\text{sf}}$ dengan metode \Caratheodory tepat merupakan versi
c.l.d.\ $\tilde\mu$ dari $\mu$.

\quad{(v)} Tunjukkan bahwa $\mu^*_{\text{sf}}=\tilde\mu^*$, sehingga jika
$\mu$ lengkap dan ditentukan secara lokal, maka
$\mu^*_{\text{sf}}=\mu^*$.
%213F, 213Xc, 213Xf

\sqheader 213Xh Misalkan $(X,\Sigma,\mu)$ suatu ruang ukur dengan himpunan
terabaikan yang ditentukan secara lokal.   Tunjukkan bahwa ruang ini semihingga.
%213I

\sqheader 213Xi Misalkan $(X,\Sigma,\mu)$ suatu ruang ukur, $\hat\mu$
kelengkapan $\mu$, $\tilde\mu$ versi c.l.d.\ dari $\mu$, dan $\check\mu$
ukuran yang ditentukan dengan metode \Caratheodory dari $\mu^*$.   Tunjukkan
bahwa pernyataan-pernyataan berikut ekuivalen:
(i) $\mu$ mempunyai himpunan terabaikan yang ditentukan secara lokal;
(ii) $\mu$ dan $\tilde\mu$ memiliki
himpunan-himpunan terabaikan yang sama;
(iii) $\check\mu=\tilde\mu$;
(iv) $\hat\mu$ dan $\tilde\mu$ mempunyai himpunan berukuran berhingga yang sama;
(v) $\mu$ dan $\tilde\mu$ mempunyai fungsi yang dapat diintegralkan yang sama;
(vi) $\tilde\mu^*=\mu^*$;
(vii) ukuran luar $\mu^*_{\text{sf}}$ dari 213Xg sama dengan $\mu^*$.
%213J 213Xi

\sqheader 213Xj Misalkan $(X,\Sigma,\mu)$ suatu ruang ukur yang dapat
dilokalkan secara ketat, dengan dekomposisi $\langle X_i\rangle_{i\in I}$.
Tunjukkan bahwa
$\mu^*A=\sum_{i\in I}\mu^*(A\cap X_i)$ untuk setiap $A\subseteq X$.
%213L

\sqheader 213Xk Misalkan $(X,\Sigma,\mu)$ suatu ruang ukur lengkap yang
ditentukan secara lokal, dan misalkan $A\subseteq X$ sedemikian sehingga
$\min(\mu^*(E\cap A),\mu^*(E\setminus A))<\mu E$ setiap kali $E\in\Sigma$
dan $0<\mu E<\infty$.
Tunjukkan bahwa $A\in\Sigma$.   \Hint{diberikan $\mu F<\infty$, tinjau
irisan $E$ dari selubung terukur bagi $F\cap A$ dan $F\setminus A$ untuk
memperoleh $\mu^*(F\cap A)+\mu^*(F\setminus A)=\mu F$.}
%213M

\spheader 213Xl Kita katakan bahwa suatu ruang ukur $(X,\Sigma,\mu)$
mempunyai {\bf sifat selubung terukur} jika setiap subhimpunan dari $X$ mempunyai
selubung terukur.   (i) Tunjukkan bahwa ruang semihingga yang mempunyai sifat
selubung terukur juga mempunyai himpunan terabaikan yang ditentukan secara lokal.
(ii) Tunjukkan bahwa ruang semihingga lengkap yang mempunyai sifat selubung
terukur ditentukan secara lokal.
%213M

\spheader 213Xm Misalkan $(X,\Sigma,\mu)$ suatu ruang ukur semihingga, dan
andaikan ruang ini memenuhi kesimpulan Teorema 213N.   Tunjukkan bahwa ruang
tersebut dapat dilokalkan.   \Hint{diberikan $\Cal E\subseteq\Sigma$, tetapkan
$\Cal F=\{F:F\in\Sigma,\,E\cap F$ terabaikan untuk setiap
$E\in\Cal E\}$.
Misalkan $\Phi$ himpunan fungsi $f$ dari subhimpunan $X$ ke $\{0,1\}$
sedemikian sehingga $f^{-1}[\{1\}]\in \Cal E$ dan
$f^{-1}[\{0\}]\in\Cal F$.}
%213N

\spheader 213Xn Misalkan $(X,\Sigma,\mu)$ suatu ruang ukur.   Tunjukkan bahwa
versi c.l.d.\ tersebut dapat dilokalkan secara ketat jika dan hanya jika terdapat
keluarga saling lepas
$\Cal E\subseteq\Sigma$ sehingga $\mu E<\infty$ untuk setiap $E\in\Cal E$
dan setiap kali $F\in\Sigma$ dan $0<\mu F<\infty$ terdapat $E\in\Cal E$
sedemikian rupa sehingga $\mu(E\cap F)>0$.
%213O

\spheader 213Xo Tunjukkan bahwa versi c.l.d.\ dari sebarang ukuran bertumpu
titik juga bertumpu titik.

\vleader{48pt}{213Y}{Latihan lanjutan (a)}
%\spheader 213Ya
Misalkan $(X,\Sigma,\mu)$ suatu ruang ukur.   Tunjukkan bahwa $\mu$
semihingga jika dan hanya jika terdapat keluarga $\Cal E\subseteq\Sigma$
sedemikian sehingga
$\mu E<\infty$ untuk setiap $E\in\Cal E$ dan
$\mu F=\sum_{E\in\Cal E}\mu(F\cap E)$ untuk setiap $F\in\Sigma$.
\Hint{ambil $\Cal E$ maksimal dengan syarat bahwa irisan sebarang dua
anggotanya terabaikan.}
%213A  would be useful in 552M, (i) => (iii)

\spheader 213Yb
Tetapkan $X=\Bbb N$, dan untuk $A\subseteq X$ tetapkan

\Centerline{$\theta A=\sqrt{\#(A)}$ jika $A$ berhingga, $\infty$ jika $A$
tak berhingga.}

\noindent Tunjukkan bahwa $\theta$ adalah ukuran luar pada $X$, bahwa
$\theta A=\sup\{\theta B:B\subseteq A,\,\theta B<\infty\}$ untuk setiap
$A\subseteq X$, tetapi ukuran $\mu$ yang ditentukan dari $\theta$ dengan
metode \Caratheodory tidak semihingga.   Tunjukkan bahwa jika $\check\mu$
merupakan ukuran yang ditentukan dengan metode \Caratheodory dari $\mu^*$
(213Xa), maka
$\check\mu\ne\mu$.
%213C

\spheader 213Yc Tetapkan $X=[0,1]\times\{0,1\}$, dan misalkan $\Sigma$
keluarga subhimpunan $E$ dari $X$ sedemikian sehingga

\Centerline{$\{x:x\in[0,1],\,E[\{x\}]\ne\emptyset,\,E[\{x\}]\ne\{0,1\}\}
$}

\noindent terhitung, dengan $E[\{x\}]=\{y:(x,y)\in E\}$ untuk setiap
$x\in[0,1]$.   Tunjukkan bahwa $\Sigma$ merupakan aljabar-sigma
subhimpunan $X$.   Untuk $E\in\Sigma$, tetapkan
$\mu E=\#(\{x:(x,1)\in E\})$ jika bilangan ini berhingga, dan $\infty$ jika
tidak.   Tunjukkan bahwa $\mu$ merupakan ukuran lengkap semihingga.   Tunjukkan
bahwa ukuran $\check\mu$ yang ditentukan dari $\mu^*$ dengan metode
\Caratheodory (213Xa) tidak semihingga.   Tunjukkan bahwa domain versi c.l.d.\
dari $\mu$ adalah seluruh $\Cal PX$.
%213Xa, 213F

\spheader 213Yd Tetapkan $X=\Bbb N$, dan untuk $A\subseteq X$ tetapkan

\Centerline{$\phi A=\#(A)^2$ jika $A$ berhingga, $\infty$ jika $A$
tak berhingga.}

\noindent Tunjukkan bahwa $\phi$ memenuhi syarat-syarat 212Ya, tetapi ukuran
yang ditentukan dari $\phi$ dengan metode tersebut tidak semihingga.
%213Xe, 213F

\spheader 213Ye Misalkan $(X,\Sigma,\mu)$ suatu ruang ukur lengkap yang
ditentukan secara lokal.   Misalkan $D\subseteq X$ dan $f:D\to\Bbb R$ suatu
fungsi.   Tunjukkan bahwa pernyataan-pernyataan berikut ekuivalen:
(i) $f$ adalah terukur;  (ii)

\Centerline{$\mu^*\{x:x\in D\cap E,\,f(x)\le
a\}+\mu^*\{x:x\in D\cap E,\,f(x)\ge b\}\le\mu E$}

\noindent setiap kali $a<b$ dalam $\Bbb R$, $E\in\Sigma$, dan
$\mu E<\infty$; (iii)

\Centerline{$\min(\mu^*\{x:x\in D\cap E,\,f(x)\le
a\},\mu^*\{x:x\in D\cap E,\,f(x)\ge b\})<\mu E$}

\noindent setiap kali $a<b$ dalam $\Bbb R$ dan $0<\mu E<\infty$.  ({\it
Petunjuk\/}:
untuk (iii)$\Rightarrow$(i), tunjukkan bahwa jika $E\subseteq X$ maka

\Centerline{$\mu^*\{x:x\in D\cap E,\,f(x)>a\}
=\sup_{b>a}\mu^*\{x:x\in D\cap E,\,f(x)\ge b\}$,}

\noindent dan gunakan 213Xk di atas.)
%213Xk, 213M

\spheader 213Yf Misalkan $(X,\Sigma,\mu)$ suatu ruang ukur lengkap yang
ditentukan secara lokal, dan misalkan $\Cal E\subseteq\Sigma$ sedemikian sehingga
$\mu E<\infty$ untuk setiap $E\in\Cal E$ dan setiap kali $F\in\Sigma$ dan
$\mu F<\infty$, terdapat subkeluarga terhitung $\Cal E_0\subseteq\Cal E$
sedemikian sehingga
$F\setminus\bigcup\Cal E_0$ dan $F\cap\bigcup(\Cal E\setminus\Cal E_0)$
keduanya terabaikan.   Tunjukkan bahwa $(X,\Sigma,\mu)$ dapat dilokalkan secara ketat.
%213O
}%end of exercises

\endnotes{
\Notesheader{213} Menurut saya, dapat dikatakan bahwa apabila definisi
`ruang ukur' ditulis ulang untuk mengecualikan semua ruang yang tidak
semihingga, teori ini tidak akan kehilangan sesuatu yang berarti.   Akan
tetapi, ada alasan kuat untuk tidak mengambil langkah sedrastis itu. Salah
satunya ialah bahwa langkah tersebut akan membingungkan semua orang (jika
Anda berkata kepada khalayak yang belum dipersiapkan, `misalkan
$(X,\Sigma,\mu)$ suatu ruang ukur', ada bahaya bahwa sebagian orang mengira
yang Anda maksud ialah `ruang ukur $\sigma$-hingga', tetapi sangat sedikit
yang akan mengira bahwa yang dimaksud ialah `ruang ukur semihingga').
Bagaimanapun, inti teori ukuran ialah membedakan himpunan berdasarkan
ukurannya. Jika setiap subhimpunan dari $E$ tidak terukur, terabaikan, atau
berukuran tak berhingga, klasifikasi tersebut terlalu kasar untuk menopang
sebagian besar gagasan lazim, terutama integrasi biasa.

Kita katakan bahwa suatu himpunan terukur $E$ {\bf murni tak berhingga} jika
$E$ sendiri dan setiap subhimpunan terukurnya yang tak terabaikan berukuran
tak berhingga.   Menurut definisi integral yang saya pilih dalam Jilid 1,
setiap fungsi sederhana, dan karena itu setiap fungsi yang dapat diintegralkan,
harus nol hampir di mana-mana pada $E$.   Ini berarti bahwa seluruh teori
integrasi sama sekali mengabaikan $E$.   Dari definisi `versi c.l.d.'
(213D-213E), terlihat bahwa versi c.l.d.\ dari ukuran tersebut menjadikan
$E$ terabaikan;  demikian pula `versi semihingga' yang diuraikan dalam 213Xc.
Namun, perubahan-perubahan ini tidak memengaruhi himpunan berukuran berhingga,
dan akibatnya mempertahankan keterintegralan fungsi beserta nilai integralnya.

Alasan terkuat yang sejauh ini kita jumpai untuk ikut mempertimbangkan ruang
yang tidak semihingga ialah bahwa metode \Caratheodory tidak selalu
menghasilkan ruang semihingga.   (Saya memberikan contoh dalam 213Yb-213Yc;
contoh yang lebih penting adalah ukuran Hausdorff dalam \S\S264-265 di bawah.)
Dalam praktik, langkah yang tepat sering kali ialah mengambil versi c.l.d.\
dari ukuran yang dihasilkan oleh konstruksi \Caratheodory.

Filosofi umum yang masuk akal dalam teori ukuran ialah bahwa kita hendak
mengukur sebanyak mungkin himpunan dan mengintegralkan sebanyak mungkin fungsi
secara kanonik -- maksudnya, tanpa membuat pilihan yang terang-terangan
sewenang-wenang tentang nilai ukuran atau integral.   Penggantian ukuran
$\mu$ dengan versi c.l.d.\ yang bersesuaian, $\tilde\mu$, kurang lebih merupakan batas
terjauh yang dapat kita capai pada ruang ukur sebarang apabila tidak ada
struktur lain yang membimbing pilihan kita.

Perhatikan bahwa $\tilde\mu$ tidak sedekat $\mu$ sebagaimana kelengkapan
$\hat\mu$ dari $\mu$;  hal ini wajar, sebab jika $E\in\Sigma$ murni tak berhingga terhadap
$\mu$, kita harus memilih antara menetapkan $\tilde\mu E=0\ne\mu E$ dan
mencari suatu cara untuk memasukkan banyak himpunan berukuran berhingga ke
dalam $E$. Jika $E$ himpunan satu anggota, pilihan kedua bahkan mustahil; dan
dalam keadaan apa pun proses itu akan bersifat sewenang-wenang.
Meskipun fungsi yang dapat diintegralkan terhadap $\tilde\mu$ tidak selalu
sama dengan yang dapat diintegralkan terhadap $\mu$ (karena $\tilde\mu$
menjadikan himpunan murni tak berhingga terabaikan, sehingga fungsi indikatornya
dapat diintegralkan), keduanya `hampir' sama. Tepatnya, setiap fungsi yang dapat
diintegralkan terhadap $\tilde\mu$ dapat diubah menjadi fungsi yang dapat
diintegralkan terhadap $\mu$ dengan menyesuaikan nilainya pada suatu himpunan
yang terabaikan terhadap $\tilde\mu$.   Ini tentu bersesuaian dengan fakta
bahwa setiap himpunan berukuran berhingga terhadap $\tilde\mu$ merupakan beda
simetris antara suatu himpunan berukuran berhingga terhadap $\mu$ dan suatu
himpunan yang terabaikan terhadap $\tilde\mu$.   Untuk himpunan berukuran tak
berhingga, hal ini dapat gagal kecuali $\mu$ dapat dilokalkan (213Hb, 213Xb).

Jika $(X,\Sigma,\mu)$ semihingga, dapat dilokalkan, atau dapat dilokalkan
secara ketat, maka tentu ruang tersebut lebih dekat, dalam arti yang
bersesuaian, dengan $(X,\tilde\Sigma,\tilde\mu)$, sebagaimana dirinci dalam
213Ha-c.

Patut dicatat bahwa ukuran $\check\mu$ yang diperoleh dengan metode
\Caratheodory langsung dari ukuran luar $\mu^*$ yang ditentukan oleh $\mu$ dapat
tidak semihingga, sekalipun $\mu$ semihingga (213Yc). Namun, modifikasi
sederhana atas $\mu^*$ (213Xg) menghasilkan versi c.l.d.\ $\tilde\mu$ dari
$\mu$, yang juga dapat diperoleh dari suatu ukuran dalam yang sesuai (213Xe).
Tentu saja $\check\mu$ juga berkaitan dengan $\tilde\mu$ melalui cara-cara
lain;  lihat 213Xd.
}%end of notes

\discrpage




